Práca sa zamerala na praktickú otázku, či možno zo slovenského právneho predpisu automaticky vytvoriť znalostný graf, ktorý jazykovému modelu pri odpovedaní pomôže identifikovať vecne správne ustanovenie, a nie iba textovo podobný úryvok. Výsledky experimentov nad zákonom č. 222/2004 Z.~z. túto hypotézu podporujú. Navrhnuté riešenie, od predspracovania PDF cez trojkrokovú schémou riadenú extrakciu až po vloženie do Neo4j, dosiahlo na 21 daňových otázkach 20 správnych odpovedí. Klasický chat s~LLM dosiahol 9 z 21 správnych odpovedí a vektorový RAG nad tým istým zákonom dosiahol pri $k=10$ 6 z 21, respektíve pri $k=25$ 17 z 21 správnych odpovedí. Rozdiel sa najvýraznejšie prejavil pri otázkach, kde právne odôvodnenie vyžadovalo kombináciu viacerých paragrafov, odsekov alebo krížových odkazov. V takýchto prípadoch grafová reprezentácia umožnila sledovať explicitné väzby, ktoré samotné vektorové vyhľadávanie zachytáva iba nepriamo.

Kvantitatívna evaluácia zároveň ukázala limity navrhnutého riešenia. Slabé miesto nespočíva primárne v dodržiavaní schémy, keďže zhoda s definovanou ontológiou sa pri všetkých konfiguráciách pohybovala nad hodnotou 0{,}92, ale v sémantickej presnosti vzťahov medzi entitami. Najlepšie mikro F1 skóre pre uzly dosiahlo 0{,}824, zatiaľ čo pri vzťahoch dosiahol najlepší výsledok 0{,}460. To potvrdzuje, že identifikácia entít je podstatne stabilnejšia než ich správne prepojenie. Najnižšiu normalizovanú grafovú editovaciu vzdialenosť dosiahli konfigurácie GPT~5.5 s vysokým úsilím v zero-shot aj few-shot režime (0{,}359), pričom zlepšenie bolo výraznejšie pri zmene modelu a režimu promptovania než pri samotnom zvyšovaní rozumového úsilia. Pri odpovedaní na otázky dosiahol GraphRAG pokrytie podporujúcich ustanovení 74{,}6\,\%, pričom najčastejšie chýbali doplnkové odseky, ustanovenia k evidencii, súhrnným výkazom alebo prílohám.

Z uvedených výsledkov vyplývajú aj ďalšie smery vývoja. Pri spracovaní väčšieho množstva právnych predpisov je zásadným obmedzením cena komerčných API, keďže spracovanie jedného približne 150-stranového dokumentu bolo v rozmedzí 8 až 15\,€. Pre širšie nasadenie preto dáva zmysel uvažovať o lokálnych modeloch, fine-tuningu na slovenskej právnej doméne, prípadne o špecializovanom extraktore doplnenom validačným krokom na kontrolu konzistentnosti entít a vzťahov. Aj v súčasnej podobe však výsledky ukazujú, že znalostný graf nie je nad právnym textom iba doplnkovou vrstvou. Pri dokumentoch s hustou sieťou odkazov poskytuje štruktúrovaný spôsob, ako odpoveď nielen vygenerovať, ale aj oprieť o konkrétne ustanovenia zákona.
